\justifying
\NSFCBodyText

\subsubsection{已具备的实验条件}

新疆大学计算机科学与技术学科具有多语言多模态认知计算与内容安全、信号检测与智能处理、大数据分析与智能控制等研究方向，可为多模态建模、传感信号处理和实验评价提供学科支撑。学院人工智能算力中心现有材料记载拥有 GPU/NPU 训练及推理服务器或工作站 130 余台、各型加速卡 600 余块，算力约为 43 PFlops@FP32/120 PFlops@FP16；教育部丝路多语言认知计算国际合作联合实验室另配置 A100 八卡服务器租赁服务。上述条件能够满足多模态编码器、时序模型、重复实验和聚类 bootstrap 的离线计算需求，其在项目执行期内的具体配额与使用流程需进一步确认。

申请人所在团队具有多模态视觉、时间序列与软件研发基础，能够完成数据字典、模型实现、对照消融和弱网轨迹回放。现有研究场所和通用计算条件可支持受控数据管理与算法实验；公开数据及已核验文献可用于基线复现和实验协议设计。

\subsubsection{尚缺少的实验条件及解决途径}

目前材料尚未证明已具备固定端侧计算设备、惯性传感器与包装图像采集装置、可追踪的传感器故障注入条件、真实弱网轨迹及功耗测量设备。拟优先选择可重复配置的端侧硬件与传感器开展受控实验，先以实际传输字节和 P95 告警时延作为主资源指标；功耗设备未及时落实时，能耗仅作次要估计，待设备到位后补充实测。具体设备型号、采购/借用方式和时间表需由申请人确认。

真实跨境物流线路、合作单位、货物级永久标识、异常终点和数据授权亦未在现有材料中落实。项目将把受控四格析因与网络轨迹回放作为两个核心问题的确认性验证基础；只有合作、授权和客观终点齐备后，才开展小规模真实运输外部验证。若合作条件未按期形成，则不把受控结果外推为现场结论，也不增加跨阶段风险传播任务。

\subsubsection{利用研究基地的计划与落实情况}

项目可依托教育部丝路多语言认知计算国际合作联合实验室、鹏城实验室新疆网络节点、新疆信号检测与处理重点实验室、新疆多语种信息技术重点实验室及新疆多模态智能处理与信息安全工程技术研究中心等平台开展多模态、网络与信号处理研究。计划依托人工智能算力中心完成模型训练、基线复现和统计重复，依托相关实验室开展传感数据与弱网回放验证。各平台是否已明确同意、本项目可使用的设备和配额以及执行期覆盖情况，需在正式申请前取得可核验口径；未落实的平台不得表述为专属条件。
